% Рисунок к определению угловой (крайней) точки.
% Слева: вершины многогранного множества — крайние точки.
% Справа: точка внутри ребра и внутренняя точка крайними не являются.
\begin{figure}[htbp]
\centering
\begin{tikzpicture}[scale=1.0, >={Stealth[length=2mm]}]

% --- Левая панель: многогранное множество и его вершины
\begin{scope}
  \coordinate (A) at (0,0);
  \coordinate (B) at (3,0.3);
  \coordinate (C) at (3.6,2.2);
  \coordinate (D) at (1.6,3.1);
  \coordinate (E) at (-0.3,1.8);
  \fill[blue!10] (A) -- (B) -- (C) -- (D) -- (E) -- cycle;
  \draw[blue!55!black, thick] (A) -- (B) -- (C) -- (D) -- (E) -- cycle;
  \foreach \p in {A,B,C,D,E}
    \fill[blue!55!black] (\p) circle (2pt);
  \node[font=\footnotesize] at (1.7,1.4) {$S$};
  \node[font=\footnotesize, anchor=north] at (1.65,-0.45)
    {вершины — крайние точки};
\end{scope}

% --- Правая панель: точка на ребре и внутренняя точка
\begin{scope}[xshift=6.3cm]
  \coordinate (A) at (0,0);
  \coordinate (B) at (3,0.3);
  \coordinate (C) at (3.6,2.2);
  \coordinate (D) at (1.6,3.1);
  \coordinate (E) at (-0.3,1.8);
  \fill[blue!10] (A) -- (B) -- (C) -- (D) -- (E) -- cycle;
  \draw[blue!55!black, thick] (A) -- (B) -- (C) -- (D) -- (E) -- cycle;

  % точка внутри ребра BC
  \coordinate (P)  at ($(B)!0.5!(C)$);
  \coordinate (P1) at ($(B)!0.15!(C)$);
  \coordinate (P2) at ($(B)!0.85!(C)$);
  \draw[orange!75!black, very thick] (P1) -- (P2);
  \fill[orange!75!black] (P1) circle (1.6pt) (P2) circle (1.6pt);
  \fill[orange!75!black] (P) circle (2.2pt);
  \node[font=\footnotesize, anchor=west, orange!60!black]
    at ($(P)+(0.25,0.1)$) {$x$};
  \node[font=\scriptsize, anchor=west] at ($(P1)+(0.2,-0.18)$) {$x_1$};
  \node[font=\scriptsize, anchor=west] at ($(P2)+(0.12,0.12)$) {$x_2$};

  % внутренняя точка
  \coordinate (Q)  at (1.5,1.3);
  \coordinate (Q1) at (0.5,1.0);
  \coordinate (Q2) at (2.5,1.6);
  \draw[orange!75!black, very thick] (Q1) -- (Q2);
  \fill[orange!75!black] (Q1) circle (1.6pt) (Q2) circle (1.6pt)
                         (Q) circle (2.2pt);

  \node[font=\footnotesize, anchor=north, text width=5cm, align=center]
    at (1.65,-0.45) {$x = \tfrac12(x_1+x_2)$, $x_1 \ne x_2$: не крайняя точка};
\end{scope}

\end{tikzpicture}
\caption{Крайние точки выпуклого множества. Слева отмечены вершины — их нельзя
представить как середину отрезка, лежащего в $S$. Справа видно, что ни внутренняя
точка, ни точка внутри ребра крайними не являются: через каждую из них проходит
отрезок множества, для которого она — середина.}
\label{fig:extreme-point}
\end{figure}
